% =============================================================================
% Dossier de conception — Éditeur de documents partagé
%
% Structure : un fichier maître (celui-ci) et quatre parties incluses :
%   partie-0-introduction.tex
%   partie-a-architecture.tex
%   partie-b-conception.tex
%   partie-c-conclusion.tex
%
% Compilation :
%   pdflatex conception.tex && pdflatex conception.tex
% (deux passes : la première résout les références et la table des matières)
%
% Les figures sont lues dans ../diagrammes-conception/png/.
% =============================================================================
\documentclass[11pt,a4paper]{article}

\usepackage[utf8]{inputenc}
\usepackage[T1]{fontenc}
\usepackage[french]{babel}
\usepackage{lmodern}
\usepackage{microtype}

\usepackage[a4paper,margin=2.4cm,headheight=15pt]{geometry}
\usepackage{graphicx}
\usepackage{float}
\usepackage{booktabs}
\usepackage{longtable}
\usepackage{array}
\usepackage{enumitem}
\usepackage{xcolor}
\usepackage{caption}
\usepackage{titlesec}
\usepackage{fancyhdr}
\usepackage[colorlinks=true,linkcolor=black,urlcolor=blue!60!black,citecolor=black]{hyperref}
\usepackage{url}

% --- Palette sobre, lisible en noir et blanc --------------------------------
\definecolor{accent}{HTML}{1B3A6B}
\definecolor{attenue}{HTML}{555555}
\definecolor{encadre}{HTML}{F7FBFF}

% --- Titres -----------------------------------------------------------------
\titleformat{\section}{\Large\bfseries\color{accent}}{\thesection}{0.7em}{}
\titleformat{\subsection}{\large\bfseries\color{accent}}{\thesubsection}{0.6em}{}
\titleformat{\subsubsection}{\normalsize\bfseries}{\thesubsubsection}{0.5em}{}
\titlespacing*{\section}{0pt}{2.2ex plus 1ex}{1.2ex}
\titlespacing*{\subsection}{0pt}{1.8ex plus 1ex}{1ex}

% --- En-têtes ---------------------------------------------------------------
\pagestyle{fancy}
\fancyhf{}
\fancyhead[L]{\small\color{attenue}Éditeur de documents partagé}
\fancyhead[R]{\small\color{attenue}\thepage}
\renewcommand{\headrulewidth}{0.4pt}

% --- Légendes ---------------------------------------------------------------
\captionsetup{
  font=small,
  labelfont=bf,
  justification=raggedright,
  singlelinecheck=false,
  skip=6pt
}

% --- Raccourcis -------------------------------------------------------------
% Insère une figure depuis le dossier des rendus PlantUML.
%   \illus{02-deploiement}{0.95}{Titre de la légende}{label}
%
% La commande ne s'appelle pas \figure : `babel-french` définit déjà ce nom
% (pour franciser les légendes), et la redéfinir provoque un arrêt net.
\newcommand{\illus}[4]{%
  \begin{figure}[H]
    \centering
    \includegraphics[width=#2\linewidth]{../diagrammes-conception/png/#1.png}
    \caption{#3}
    \label{#4}
  \end{figure}%
}

% Identifiants techniques : chasse fixe, taille courante.
\newcommand{\code}[1]{\texttt{#1}}

% Encadré « hypothèse » : ce qui relève d'une interprétation, et non du sujet.
\newcommand{\hypothese}[1]{%
  \par\medskip
  \noindent\colorbox{encadre}{%
    \begin{minipage}{\dimexpr\linewidth-2\fboxsep}
      \small\textbf{Hypothèse d'interprétation.} #1
    \end{minipage}%
  }
  \par\medskip
}

\setlength{\parskip}{0.35em}
\setlength{\parindent}{0pt}

% Les tableaux sont larges : l'intercolonne par défaut (6 pt de chaque côté de
% chaque colonne) fait déborder les tableaux à trois colonnes de 5 mm dans la
% marge. On la ramène à 4 pt, et les largeurs de colonnes sont exprimées en
% fraction de \linewidth plutôt qu'en centimètres : le tableau reste juste même
% si les marges changent.
\setlength{\tabcolsep}{4pt}
\renewcommand{\arraystretch}{1.12}

\hypersetup{
  pdftitle={Dossier de conception — Éditeur de documents partagé},
  pdfauthor={Djusse Tameno Christian},
  pdfsubject={Projet Cloud — ENSPY}
}

% =============================================================================
\begin{document}

% --- Page de titre ----------------------------------------------------------
\begin{titlepage}
  \centering
  \vspace*{2.5cm}

  {\color{accent}\rule{\linewidth}{1.2pt}}\\[0.6cm]
  {\Huge\bfseries Éditeur de documents partagé\par}
  \vspace{0.4cm}
  {\Large Dossier de conception\par}
  \vspace{0.3cm}
  {\color{accent}\rule{\linewidth}{1.2pt}}\\[1.2cm]

  {\large Écriture collaborative en temps réel\\
  et distribution de documents à quota strict\\[0.3cm]
  sur une infrastructure AWS\par}

  \vspace{2.2cm}

  % La colonne de droite est une colonne de **paragraphe** : dans une colonne
  % `r`, ni `\url` ni `\texttt` ne peuvent se couper, et l'adresse du dépôt
  % débordait de la marge de 7,5 pt. Une largeur fixe autorise la césure.
  \begin{tabular}{@{}rl@{}}
    \textbf{Auteur} & Djusse Tameno Christian \\[0.2cm]
    \textbf{Filière} & 5\textsuperscript{e} année Génie informatique, ENSPY \\[0.2cm]
    \textbf{Cours} & Cloud — projet d'évaluation individuel \\[0.2cm]
    \textbf{Région AWS} & \code{eu-north-1} (Stockholm) \\[0.2cm]
    \textbf{Dépôt} & \href{https://github.com/Djusse/Editeur_Partag-}%
                       {\texttt{github.com/Djusse/}}\texttt{/}%
                       \href{https://github.com/Djusse/Editeur_Partag-}%
                       {\texttt{Editeur\_Partag-}} \\
  \end{tabular}

  \vfill

  {\small\color{attenue}
    Ce dossier décrit la conception telle qu'elle est \emph{implémentée}, et non
    telle qu'elle avait été prévue. Les écarts sont signalés, avec leur raison.\\
    Chaque affirmation sur les contraintes renvoie à un test ou à un scénario
    de charge exécutable.\par}
\end{titlepage}

% --- Tables -----------------------------------------------------------------
\tableofcontents
\newpage

\listoffigures
\listoftables
\newpage

% --- Parties ----------------------------------------------------------------
% =============================================================================
% Partie 0 — Introduction
% =============================================================================
\section{Introduction}

\subsection{Objet du document}

Ce dossier décrit la conception d'un éditeur de documents partagé déployé sur
AWS. L'application poursuit deux buts que l'énoncé oppose volontairement :

\begin{itemize}[leftmargin=1.4em,itemsep=2pt]
  \item plusieurs centaines de personnes écrivent \emph{en même temps} dans le
        même texte, et aucune frappe acceptée ne doit être perdue ;
  \item un lien public distribue le document en PDF à \emph{exactement} $N$
        personnes, jamais $N+1$, chacune recevant le document tel qu'il était
        au moment où elle a réservé son exemplaire.
\end{itemize}

Le premier objectif demande de la souplesse : des écritures concurrentes, un
historique, de la reprise après coupure. Le second demande de la rigidité : un
compteur exact, une transaction, une réponse unique. Tenir les deux dans la même
application est l'essentiel du travail de conception.

\subsection{Ce que ce dossier est, et n'est pas}

Ce dossier décrit la conception \emph{réellement implémentée}, vérifiée contre
le code source, et non une intention de conception. Trois conséquences :

\begin{enumerate}[leftmargin=1.4em,itemsep=2pt]
  \item là où l'implémentation s'écarte du plan initial, l'écart est signalé avec
        sa raison — ces écarts sont souvent la partie la plus instructive ;
  \item chaque affirmation sur une contrainte renvoie à un test unitaire, un
        script d'intégration ou un scénario de charge, tous exécutables ;
  \item les limites sont énoncées sans être atténuées. Une limite annoncée vaut
        mieux qu'une limite découverte par le lecteur.
\end{enumerate}

\subsection{Plan}

La \textbf{partie A} traite de l'architecture : contexte, déploiement, décisions
et alternatives écartées, flux internes, sécurité, infrastructure, limites.

La \textbf{partie B} traite de la conception applicative : cas d'utilisation,
organisation des paquets, modèle du domaine, états, séquences système, puis une
lecture par contrainte, de C1 à C6.

La \textbf{partie C} conclut : ce qui est prouvé, ce qui reste à mesurer, et ce
que ce projet a réellement appris.

% -----------------------------------------------------------------------------
\section{Contexte, objectif et contraintes}

\subsection{Contexte}

Le sujet impose de démontrer la maîtrise des services AWS au travers d'une
application non triviale, dont chaque choix doit pouvoir être justifié. Il
précise que « une application simple, mais dont chaque choix est compris et
justifié, vaut mieux qu'une application complexe que vous ne savez pas
expliquer ». Ce dossier suit cette consigne : peu de services, mais chacun
employé pour ce qu'il fait de mieux, et chaque fois avec l'alternative écartée.

Deux points du sujet sont éliminatoires : les contraintes C1 et C2. Une
application rapide qui les viole est non fonctionnelle, quel que soit son
polissage.

\subsection{Les six contraintes}

\begin{table}[H]
\centering
\small
\begin{tabular}{@{}p{0.6cm}p{5.3cm}p{7.3cm}@{}}
\toprule
\textbf{\#} & \textbf{Contrainte de l'énoncé} & \textbf{Réponse retenue} \\
\midrule
C1 & Exactement $N$ téléchargements, jamais $N+1$
   & transaction DynamoDB : décrément de \code{remaining} et création du ticket
     dans la même écriture conditionnelle. Une garde en mémoire allège la base
     sous le pic, mais n'est jamais l'autorité \\
\addlinespace
C2 & Aucune modification acceptée perdue ; convergence
   & CRDT Yjs avec \emph{un seul écrivain} par document ; une modification n'est
     acquittée qu'\emph{après} écriture durable dans le journal \\
\addlinespace
C3 & Document tel qu'il est à $T_1$, complet et cohérent
   & le claim fige un numéro de séquence $T_1$ ; le PDF est construit depuis le
     journal jusqu'à ce numéro, une seule fois par version, de façon asynchrone \\
\addlinespace
C4 & Jamais plus que son crédit, même sur plusieurs appareils
   & compteur DynamoDB par (e-mail, jour) débité dans la même transaction que
     l'écriture au journal ; comptage effectué côté serveur, sur le binaire Yjs \\
\addlinespace
C5 & 15\,000 clics (80\,\% en 60\,s) pendant que 300 rédacteurs écrivent
   & micro-cache nginx sur l'état du lien, garde en mémoire sur les
     réservations, diffusion par lots de 250\,ms, instance dimensionnée par la
     mesure. Trois états explicites : prêt, patienter, épuisé \\
\addlinespace
C6 & Dernier exemplaire : visiteurs prévenus vite, plus de travail inutile
   & état épuisé monotone : réponse immédiate, aucun appel à la base, mis en
     cache longuement par nginx \\
\bottomrule
\end{tabular}
\caption{Les six contraintes et la réponse architecturale retenue.}
\label{tab:contraintes}
\end{table}

\subsection{Hypothèses d'interprétation}

Le sujet laisse plusieurs points ouverts. Les choix ci-dessous ne sont pas
confirmés par l'enseignant : ils sont donc explicités, pour que le lecteur
puisse contester l'interprétation plutôt que le résultat.

\begin{table}[H]
\centering
\small
\begin{tabular}{@{}p{5.4cm}p{10.2cm}@{}}
\toprule
\textbf{Sujet} & \textbf{Interprétation retenue} \\
\midrule
Consommation d'un exemplaire & au clic qui réserve le téléchargement, et non au
  téléchargement effectif du fichier \\
\addlinespace
Instant $T_1$ & lecture du dernier numéro de séquence au début du claim \\
\addlinespace
Crédit & 10\,000 caractères par jour ; insertions \emph{et} suppressions
  comptées ; restauration de version gratuite \\
\addlinespace
Recharge du crédit & minuit, heure de Douala — et non minuit UTC (voir B.6) \\
\addlinespace
Blocage & 4 refus pour crédit épuisé dans la journée bloquent l'écriture
  jusqu'à la recharge ; la lecture et le téléchargement restent possibles \\
\addlinespace
Pic d'écriture & 300 rédacteurs dans le \emph{même} document : c'est le pire cas,
  retenu comme cible ; le palier réellement atteint est celui qui est rapporté \\
\addlinespace
Identité & nom et e-mail, sans mot de passe. Le sujet l'autorise ; l'e-mail est
  la clé du crédit et du partage \\
\addlinespace
Livrables & français, dépôt GitHub privé, jury ajouté en collaborateur \\
\bottomrule
\end{tabular}
\caption{Hypothèses d'interprétation, à confronter à l'énoncé.}
\label{tab:hypotheses}
\end{table}

\hypothese{La recharge du crédit à minuit \emph{heure de Douala} s'écarte de
l'hypothèse initiale (minuit UTC). L'écart vient de l'usage : un utilisateur à
Douala lisait « recharge à 01:00 » et ne comprenait pas. Le décalage est
constant toute l'année, le Cameroun n'appliquant pas d'heure d'été, ce qui rend
l'arithmétique exacte sans table de fuseaux.}

\subsection{Objectif mesurable}

La conception n'est acceptable que si elle est mesurable. Quatre jalons ont
guidé les décisions :

\begin{enumerate}[leftmargin=1.4em,itemsep=2pt]
  \item $N$ réservations exactement pour un lien de $N$ exemplaires, sous
        15\,000 clics dont 80\,\% en 60 secondes, sans aucune erreur serveur ;
  \item zéro modification acceptée sans acquittement, sous écriture simultanée
        de plusieurs dizaines de rédacteurs ;
  \item convergence du contenu vérifiable : journal sans trou, reconstruction
        identique par deux chemins indépendants ;
  \item un palier de rédacteurs réellement atteint, rapporté tel quel, même s'il
        est inférieur à la cible.
\end{enumerate}

Le troisième point mérite d'être souligné : un test de charge mesure un
\emph{débit d'acquittement}, pas une convergence de contenu. Les deux preuves
sont distinctes, et le dossier les traite séparément.

% =============================================================================
% Partie A — Architecture globale
% =============================================================================
\section{Partie A — Architecture globale}

\subsection{Diagramme de contexte}

Le système a trois catégories d'utilisateurs, qui n'accèdent pas au service de la
même façon.

Le \textbf{rédacteur} s'identifie par un nom et une adresse électronique, sans
mot de passe, et reçoit un jeton de session. Il crée des documents, les partage,
les versionne et fabrique des liens de téléchargement.

Le \textbf{visiteur} est le cas intéressant : il ne s'identifie \emph{jamais}.
Le jeton du lien qu'il détient \emph{est} sa capacité d'accès. Cette décision a
une conséquence architecturale : le parcours visiteur est exposé à la racine
(\code{/d/\{token\}/\dots}), en dehors du préfixe \code{/api} qui regroupe les
routes authentifiées. La confondre avec les autres routes produit un échec
silencieux : la requête aboutit sur le fichier statique du front et renvoie un
code 200 avec du HTML.

L'\textbf{administrateur} consulte l'état du système avec une clé dédiée, dans
un en-tête, et non avec le jeton de session : c'est une capacité
d'exploitation, pas une identité applicative.

Deux éléments extérieurs complètent le tableau. \textbf{Amazon SES} envoie un
unique courriel lorsque le dernier exemplaire d'un lien est distribué.
\textbf{Le stockage objet} sert les fichiers PDF : le système ne les transporte
jamais, il présigne une adresse et laisse S3 répondre. C'est ce qui permet à
500 téléchargements de ne pas passer par l'instance applicative.

\illus{01-contexte}{0.82}{Diagramme de contexte : trois acteurs, un système,
deux services extérieurs.}{fig:contexte}

Un point de conception mérite d'être explicité ici. Le destinataire du courriel
de fin de quota n'est pas un acteur du système : il ne dispose d'aucune
interface, ne peut rien demander et ne reçoit rien d'autre. Le représenter comme
un acteur donnerait une idée fausse du périmètre.

\subsection{Architecture de déploiement AWS}

Tout le trafic entre par une \textbf{adresse IP élastique}, attachée à une
instance unique. Cette adresse est indispensable : une adresse publique
ordinaire change à chaque arrêt puis démarrage de l'instance, ce qui rendrait
l'évaluation impossible avec une URL stable.

Sur l'instance cohabitent trois rôles, séparés par des processus distincts et
non par des machines :

\begin{itemize}[leftmargin=1.4em,itemsep=2pt]
  \item \textbf{nginx} sert le front statique, met en micro-cache l'état des
        liens, et relaie vers Node les préfixes \code{/api}, \code{/d} et
        \code{/ws} ;
  \item \textbf{le processus applicatif} tient l'état Yjs des documents ouverts
        en mémoire et sert l'API REST et le WebSocket ;
  \item \textbf{le processus de travail} consomme la file de génération de PDF,
        séparément pour qu'un document lourd ne bloque pas les écritures.
\end{itemize}

\illus{02-deploiement}{0.95}{Architecture de déploiement : une région, une
instance applicative, quatre services gérés.}{fig:deploiement}

Le choix du type d'instance est une décision de dimensionnement, pas de confort.
Le type retenu est \code{m7i-flex.large} : deux processeurs virtuels et huit
gigaoctets de mémoire. Deux raisons, dans cet ordre :

\begin{enumerate}[leftmargin=1.4em,itemsep=2pt]
  \item c'est le seul type à la fois \textbf{éligible à l'offre gratuite} et doté
        de la mémoire nécessaire pour tenir des milliers de connexions
        simultanées. \code{t3.large} a été essayé d'abord : la création de pile
        échoue avec « not eligible for Free Tier » ;
  \item les instances de la famille \code{t} consomment des crédits de processeur
        et se retrouvent bridées quand ces crédits s'épuisent. En mode
        \emph{Standard}, une instance bridée s'effondre précisément pendant le
        pic qu'on cherche à mesurer. Le mode \emph{illimité} évite le
        bridage mais rend la facture imprévisible.
\end{enumerate}

\subsection{Décisions d'architecture et alternatives écartées}

Chaque décision est consignée dans un document d'architecture séparé
(\code{docs/adr/}), avec ses alternatives. Le tableau ci-dessous en donne la
synthèse et, surtout, ce qui a été \emph{écarté} : une décision sans alternative
n'est pas une décision.

\begin{table}[H]
\centering
\small
% Largeur utile = 21 cm − 2×2,4 cm de marge − 4×4 pt d'intercolonne = 15,64 cm.
% Le total ci-dessous (15,5 cm) tient avec une marge ; à 15,9 cm il débordait.
\begin{tabular}{@{}p{3.0cm}p{6.2cm}p{6.3cm}@{}}
\toprule
\textbf{Choix} & \textbf{Retenu, et pourquoi} & \textbf{Écarté, et pourquoi} \\
\midrule
Calcul
& \textbf{Instance EC2 unique} : le WebSocket et l'état en mémoire sont
  naturels, et un seul écrivain par document rend C2 simple à tenir
& \textbf{Lambda + API Gateway} : la limite de 400 exécutions simultanées du
  compte est trop serrée pour 15\,000 clics, et le WebSocket persistant
  obligerait à externaliser tout l'état \\
\addlinespace
Convergence
& \textbf{CRDT Yjs} : les modifications concurrentes fusionnent sans verrou, et
  un client hors ligne se réconcilie à la reconnexion
& \textbf{Verrou exclusif} : un seul rédacteur à la fois, ce que le sujet
  interdit ; \textbf{transformation opérationnelle} : juste, mais coûteux à
  écrire et à prouver dans le temps d'un projet \\
\addlinespace
Écriture
& \textbf{Écrivain unique et lots de 250\,ms} : une transaction par lot au lieu
  d'une par frappe
& \textbf{Écriture par frappe} : une transaction DynamoDB par caractère épuise
  le débit disponible et multiplie la facture \\
\addlinespace
Persistance
& \textbf{DynamoDB} : transactions conditionnelles, ce qui donne C1 exactement,
  et aucune administration de serveur
& \textbf{PostgreSQL} : les transactions seraient plus expressives, mais
  imposent un serveur à administrer, sauvegarder et dimensionner \\
\addlinespace
Fichiers
& \textbf{S3 privé + URL présignée de 60 s} : le fichier ne traverse pas
  l'instance
& \textbf{Flux servi par Node} : 500 PDF simultanés saturent la bande passante
  de l'instance et bloquent l'édition \\
\addlinespace
PDF
& \textbf{Génération asynchrone par file} : le visiteur reçoit « patienter » et
  revient ; un seul travail par version grâce à un verrou
& \textbf{Génération dans la requête} : une requête de plusieurs secondes sous
  15\,000 clics ferait tomber l'instance \\
\addlinespace
Pic visiteur
& \textbf{Micro-cache nginx sur l'état du lien} : une seconde de cache absorbe
  15\,000 lectures sans toucher la base
& \textbf{Interrogation directe de DynamoDB} : 15\,000 lectures par seconde la
  font échouer en limitation de débit \\
\addlinespace
Front
& \textbf{Export statique servi par nginx} : quelques fichiers, aucun serveur
  applicatif pour le rendu
& \textbf{Rendu côté serveur Next.js} : un second processus Node à
  dimensionner, sans bénéfice ici \\
\bottomrule
\end{tabular}
\caption{Décisions d'architecture et alternatives écartées.}
\label{tab:adr}
\end{table}

\subsection{Flux internes}

Deux flux méritent une figure : celui qui tient la contrainte C1, et celui qui
tient C2 avec C4. Ce sont les deux contraintes éliminatoires.

\subsubsection{Réservation d'un exemplaire}

L'ordre des vérifications est imposé par le sujet et a des conséquences
observables. Avant de décider que le lien est épuisé, le serveur cherche
\emph{d'abord} si ce visiteur possède déjà un ticket. Sans cette précaution, un
visiteur déjà servi recevrait « épuisé » après l'épuisement du lien — il perdrait
un exemplaire qu'il avait légitimement obtenu.

Ensuite, la réservation est une transaction conditionnelle : décrémenter
\code{remaining} \emph{si} il est supérieur à zéro et si le lien est actif, et
créer le ticket \emph{si} ce visiteur n'en a pas. La transition vers l'état
épuisé est une écriture séparée, elle aussi conditionnelle : seul le claim qui
amène réellement le compteur à zéro peut la poser. C'est cette séparation qui
rend l'épuisement idempotent lorsque cinquante visiteurs arrivent ensemble.

\illus{03-flux-reservation}{0.92}{Flux de réservation : transaction
conditionnelle, puis génération asynchrone du PDF.}{fig:flux-reservation}

Le numéro de séquence est figé au moment du claim, avant la transaction. Tous
les visiteurs d'une même rafale partagent donc le même fichier, et une
modification ultérieure du document ne change pas leur exemplaire.

\subsubsection{Écriture groupée}

Le flux d'écriture porte deux contraintes à la fois : aucune modification perdue
(C2) et aucun dépassement de crédit (C4). Les deux se règlent au même endroit,
parce qu'elles portent sur la même transaction.

\illus{04-flux-ecriture}{0.92}{Flux d'écriture : un lot par tick de 250\,ms, une
transaction atomique, un seul message diffusé.}{fig:flux-ecriture}

Un point de conception a coûté cher et mérite d'être rapporté. Le contrôle du
crédit ne peut pas être une condition de la transaction : DynamoDB refuse
l'arithmétique entre deux valeurs dans une expression de condition. La forme
\code{used + :chars <= :limit} est rejetée par le service. Le contrôle se fait
donc par une lecture préalable, hors transaction. Le risque de course est borné
par l'écrivain unique ; la transaction reste la garde d'atomicité pour la
séquence, le journal et les marqueurs d'idempotence.

\subsection{Sécurité et bonnes pratiques}

Les mesures ci-dessous répondent directement aux bonnes pratiques du cours.

\begin{table}[H]
\centering
\small
\begin{tabular}{@{}p{4.6cm}p{11cm}@{}}
\toprule
\textbf{Point} & \textbf{Mise en œuvre} \\
\midrule
Aucune clé d'accès dans le code
& le poste utilise une connexion fédérée ; le serveur utilise le rôle de
  l'instance, donc des identifiants temporaires renouvelés \\
\addlinespace
Aucun secret dans les données de démarrage
& les données de démarrage d'une instance sont lisibles via le service de
  métadonnées. Les deux secrets (signature des sessions, clé d'administration)
  sont dans le magasin de paramètres, créés une fois et conservés \\
\addlinespace
Moindre privilège
& le rôle de l'instance n'accède qu'aux tables nommées, au préfixe de stockage,
  à la file, et à l'envoi de courriel filtré sur l'expéditeur vérifié \\
\addlinespace
Métadonnées d'instance
& jeton obligatoire : sans session, les données de démarrage ne sont pas
  lisibles \\
\addlinespace
Stockage objet
& accès public bloqué, chiffrement au repos, versionnage, règle refusant le
  trafic non chiffré, URL présignées de 60 secondes \\
\addlinespace
Pare-feu
& port 80 ouvert ; port 22 \emph{fermé} par défaut, accès par session
  sécurisée. Une valeur permissive est explicitement traitée comme une
  fermeture, et non comme une ouverture \\
\addlinespace
Jeton WebSocket
& le navigateur ne permet pas d'en-tête d'authentification au handshake : un
  ticket à usage unique de 60 secondes évite d'exposer le jeton de session dans
  l'URL, donc dans les journaux du serveur frontal \\
\addlinespace
Validation
& adresse électronique, longueur du titre, longueur de l'identifiant
  d'opération, plafond de 64 Kio par mise à jour, corps de requête limité \\
\bottomrule
\end{tabular}
\caption{Mesures de sécurité et bonne pratique correspondante.}
\label{tab:securite}
\end{table}

\subsection{Infrastructure, déploiement et suppression}

L'infrastructure entière est décrite dans un seul modèle de pile, ce qui rend le
déploiement reproductible : une instance, une adresse élastique, un groupe de
sécurité, dix tables, un compartiment de stockage, une file avec file de
lettres mortes, un rôle d'instance et un groupe de journaux.

Le script de déploiement enchaîne six étapes, et chacune échoue franchement :

\begin{enumerate}[leftmargin=1.4em,itemsep=2pt]
  \item vérification de l'identité, de la région, et de l'\emph{éligibilité} du
        type d'instance — interrogée auprès du service, et non tenue dans une
        liste à la main ;
  \item création éventuelle de la paire de clés, puis des deux secrets ;
  \item déploiement de la pile, avec suppression préalable si elle se trouve
        dans un état d'échec antérieur, qui interdit toute mise à jour ;
  \item construction du backend puis du front, en libérant les dépendances au
        fur et à mesure pour tenir dans l'espace disque disponible ;
  \item envoi de l'archive vers le stockage, puis installation sur l'instance
        par commande à distance ;
  \item vérification que le service répond, que le parcours visiteur est
        relayé, et que le micro-cache est actif.
\end{enumerate}

La suppression mérite un script dédié, parce qu'une simple suppression de pile
échoue : le compartiment est versionné, donc il faut purger les objets
\emph{et} leurs versions ; les paramètres et la paire de clés vivent hors de la
pile et survivent à sa suppression. Le script purge, supprime, puis \emph{vérifie}
huit familles de ressources, en distinguant trois états : absent, restant, et
\emph{inconnu} — une commande de contrôle qui échoue ne doit jamais être comptée
comme « rien ne reste ».

\subsection{Limites et risques}

\begin{table}[H]
\centering
\small
\begin{tabular}{@{}p{4.4cm}p{11.2cm}@{}}
\toprule
\textbf{Limite} & \textbf{Conséquence, et raison du choix} \\
\midrule
Trafic sans chiffrement
& les jetons et les adresses circulent en clair. Un certificat exige un nom de
  domaine, absent du sujet. Limite assumée, documentée, et non dissimulée \\
\addlinespace
Instance unique
& aucune haute disponibilité : un arrêt de la machine interrompt le service.
  C'est le prix d'un écrivain unique par document, qui simplifie C2. Le
  redémarrage automatique ramène le processus, pas la machine \\
\addlinespace
Identité sans mot de passe
& quiconque connaît une adresse peut écrire en son nom. Le sujet l'autorise ;
  c'est aussi ce qui rend l'identité du visiteur inutile \\
\addlinespace
Jeton non révocable
& un jeton volé reste valable jusqu'à son expiration. Sans liste de révocation,
  il n'y a pas de contournement autre que le renouvellement \\
\addlinespace
Rattrapage plafonné
& la lecture du journal s'arrête à un nombre d'entrées fixé, sans signaler la
  troncature : un client doit boucler. Un appel unique sur un document ancien
  donnerait un journal partiel, donc une vérification de convergence fausse \\
\addlinespace
Restauration non facturée
& une restauration de version est gratuite et franchit un blocage. C'est un
  recours voulu, mais le compteur ne reflète donc pas tout le travail effectué \\
\addlinespace
Bac à sable du courriel
& tant que l'accès de production n'est pas accordé, seule l'adresse vérifiée
  reçoit la notification \\
\bottomrule
\end{tabular}
\caption{Limites connues et risques associés.}
\label{tab:limites}
\end{table}

\hypothese{L'absence de haute disponibilité est présentée comme une limite
assumée, non comme un oubli. Le cours demande explicitement de savoir et de dire
qu'une instance unique n'est pas disponible. Les pistes d'évolution — plusieurs
instances avec élection d'un écrivain par document, ou externalisation de l'état
— sont plus coûteuses que le bénéfice attendu pour un projet évalué sur la
justification des choix.}

% =============================================================================
% Partie B — Conception de l'application
% =============================================================================
\section{Partie B — Conception de l'application}

\subsection{Cas d'utilisation}

Trois acteurs, onze cas. La lecture utile n'est pas la liste mais les deux
relations entre cas, qui portent chacune une décision.

\illus{05-cas-utilisation}{0.88}{Cas d'utilisation : rédacteur, visiteur,
administrateur.}{fig:cas-utilisation}

La première relation paraît anodine et ne l'est pas : « éditer en collaboration »
\emph{étend} « consulter son crédit », et non l'inverse. Le crédit n'est jamais
interrogé avant d'écrire. Le retour du serveur — acquittement ou refus — porte
déjà le crédit restant, et c'est le refus qui met la jauge à jour. Une relation
d'inclusion aurait décrit un aller-retour qui n'existe pas.

La seconde est une absence qu'il faut assumer : il n'y a \emph{aucune} relation
entre « réserver un exemplaire » et « s'identifier ». Le visiteur ne s'identifie
jamais. Ce n'est pas un oubli dans la conception, c'est la décision qui rend le
parcours visiteur absorbable à 15\,000 clics : pas de session à créer, pas de
jeton à vérifier, une seule écriture conditionnelle.

\subsection{Diagramme de paquets}

L'organisation du code suit une règle unique : les routes dépendent du domaine,
le domaine dépend des accès aux services, jamais l'inverse. Un seul paquet parle
aux services AWS. Le domaine reçoit le stockage, la file et le courriel par
injection, ce qui permet aux tests de n'appeler aucun service.

\illus{06-packages}{0.80}{Organisation des paquets du backend.}{fig:packages}

Cette organisation n'est pas décorative. C'est elle qui a permis d'écrire
43 tests sans aucun appel réseau, et c'est elle qui a rendu possible le mode
local complet, où l'application tourne avec un stockage et une base en mémoire,
sans compte AWS.

\subsection{Diagramme de classes et modèle de données}

Le modèle tient en dix tables, dont trois portent une règle plutôt qu'une donnée.

\illus{07-classes}{0.86}{Modèle du domaine : dix entités, trois portent une
règle de concurrence ou de quota.}{fig:classes}

Trois points de conception méritent l'attention, parce qu'ils sont faciles à
modéliser faux.

Le premier concerne les \textbf{versions}. Le numéro de version \emph{est} le
numéro de séquence du document au moment de l'enregistrement. Ce n'est pas deux
champs mais un seul, utilisé deux fois. La conséquence pratique est agréable : la
numérotation ne peut pas avoir de trous, puisqu'elle suit une séquence déjà
garantie contiguë par le journal.

Le deuxième concerne l'\textbf{idempotence}. La table des marqueurs porte deux
clés par mise à jour : l'identifiant d'opération fourni par le client, et
l'empreinte du contenu. La seconde existe parce qu'un client qui se reconnecte
peut renvoyer la même modification sous un autre identifiant ; sans elle, un
renvoi débiterait le crédit deux fois.

Le troisième concerne le \textbf{cache de PDF}. Il porte un compteur de
tentatives et une date de bail. Sans eux, deux visiteurs arrivant sur un PDF
absent publieraient deux travaux de génération, et un PDF qui échoue serait
redemandé indéfiniment.

\begin{table}[H]
\centering
\small
\begin{tabular}{@{}p{3.6cm}p{5.2cm}p{6.6cm}@{}}
\toprule
\textbf{Table} & \textbf{Clé} & \textbf{Rôle de conception} \\
\midrule
\code{Documents} & \code{docId} & séquence courante, dernier instantané, longueur du
  texte \\
\code{Members} & \code{docId} + \code{email} & appartenance ; index secondaire par
  adresse pour lister les documents d'un rédacteur \\
\code{Updates} & \code{docId} + \code{seq} & journal : une entrée par lot, avec
  l'update Yjs fusionné \\
\code{Versions} & \code{docId} + \code{n} & instantané consultable ; \code{n} est
  le numéro de séquence \\
\code{Links} & \code{token} & quota : \code{remaining} et statut \\
\code{Tickets} & \code{token} + \code{visitorId} & exemplaire réservé et numéro
  figé \\
\code{PdfCache} & \code{docId\#seq} & état de génération, verrou, tentatives \\
\code{Credits} & \code{email\#jour} & compteur consommé et plafond \\
\code{Abus} & \code{email\#jour} & refus du jour et terme du blocage \\
\code{OpIds} & \code{id} & idempotence : identifiant d'opération et empreinte \\
\bottomrule
\end{tabular}
\caption{Modèle de données : dix tables, une clé et un rôle par table.}
\label{tab:tables}
\end{table}

\subsection{Diagrammes d'états-transitions}

Trois objets ont un cycle de vie qui porte une règle : le lien, le visiteur, et
le crédit.

\subsubsection{États du lien}

Un lien naît actif avec son quota, le consomme, et devient épuisé. La transition
est \emph{monotone} : un lien ne redevient jamais actif. Point important pour la
lecture du diagramme : l'état épuisé \textbf{n'est pas terminal}. Le lien reste
consultable indéfiniment et répond « épuisé », ce qui est précisément ce que
demande C6 — plus aucun travail inutile.

\illus{08a-etats-lien}{0.52}{États d'un lien de téléchargement.}{fig:etats-lien}

\subsubsection{États du visiteur}

Le visiteur suit un chemin qui dépend de deux choses indépendantes : le quota du
lien et l'état du fichier. C'est ce qui produit trois réponses distinctes là où
on attendrait une seule : prêt, patienter, épuisé.

\illus{08b-etats-visiteur}{0.62}{États d'un visiteur sur un lien.}{fig:etats-visiteur}

Un point mérite d'être souligné : une fois l'exemplaire réservé, il est
\emph{acquis}. L'URL présignée expire au bout de soixante secondes, mais le
visiteur peut en redemander une indéfiniment. Rien de ce qu'il fait ensuite ne
consomme d'exemplaire supplémentaire.

\subsubsection{États du crédit}

Le crédit a trois états, et le troisième est celui qu'on oublie : un utilisateur
peut être \emph{bloqué alors qu'il lui reste du crédit}. Le blocage est posé par
quatre refus dans la journée ; il ne s'efface pas parce que le compteur a été
remis à zéro.

\illus{08c-etats-credit}{0.62}{États du crédit d'écriture, par jour et par
adresse.}{fig:etats-credit}

La remise à zéro mérite une explication, parce qu'elle a été source de confusion.
Ce n'est pas un événement, c'est un changement de clé : le compteur est indexé
par adresse et par jour. Deux journées différentes sont deux lignes différentes,
sans qu'aucune écriture ne soit nécessaire.

\hypothese{Le « jour » est celui de Douala, pas le jour UTC. C'est un écart avec
l'hypothèse initiale du plan. Un utilisateur lisait « recharge à 01:00 » et ne
comprenait pas : minuit UTC tombe à une heure du matin sur place. Le décalage
est constant — le Cameroun n'applique pas d'heure d'été — ce qui rend
l'arithmétique exacte toute l'année. Le terme du blocage est par ailleurs
\emph{recalculé} à chaque lecture et non relu en base, ce qui a corrigé les
lignes écrites sous l'ancienne règle sans aucune migration.}

\subsection{Diagrammes de séquence système}

Trois parcours structurent l'application : réserver puis télécharger, éditer,
restaurer.

\subsubsection{Réserver et télécharger}

\illus{09-ssd-telechargement}{0.80}{Séquence système : réserver puis télécharger
un exemplaire.}{fig:ssd-telechargement}

Deux précisions que le diagramme ne peut pas montrer. Consulter l'état du lien
ne consomme rien : c'est une exigence de conception, sans quoi un aperçu de lien
dans une messagerie ou un robot d'indexation épuiserait le quota. Et
l'interrogation d'attente porte sur le \emph{fichier}, pas sur l'état du lien :
l'état du lien ne dit rien de l'avancement de la génération.

\subsubsection{Éditer un document}

\illus{10-ssd-ecriture}{0.72}{Séquence système : éditer un document en
collaboration.}{fig:ssd-ecriture}

Le comportement en cas de refus mérite d'être explicite, parce qu'il est
contre-intuitif. À la réception d'un refus de crédit, le client vide sa file
locale et reconstruit le document depuis le serveur. Ce n'est pas une
resynchronisation : le refus n'a rien modifié. Le client reconstruit parce que
les frappes déjà saisies ne seront \emph{jamais} enregistrées, et qu'il faut le
montrer sans ambiguïté plutôt que laisser un état intermédiaire trompeur.

\subsubsection{Restaurer une version}

\illus{11-ssd-restauration}{0.78}{Séquence système : restaurer une version.}{fig:ssd-restauration}

La restauration est une écriture, pas un retour en arrière. Trois propriétés en
découlent, et toutes les trois sont voulues : la compatibilité avec le CRDT, donc
la convergence des autres clients ; la conservation de l'historique, puisque le
texte écrasé est d'abord sauvegardé comme version ; la gratuité, puisqu'aucun
caractère n'est facturé et que l'opération franchit un blocage. Cette dernière
propriété est un choix : sans elle, un utilisateur ayant épuisé son crédit serait
enfermé sans recours.

\subsection{Conception par contrainte}

Cette section reprend les six contraintes une par une, et indique pour chacune
\emph{où} elle se tient dans la conception, et \emph{comment} elle est vérifiée.
C'est la lecture la plus utile du dossier : elle relie chaque exigence à un
mécanisme et à une preuve.

\subsubsection{C1 — exactement $N$ téléchargements}

\textbf{Où.} Dans une seule transaction conditionnelle sur la table des liens et
la table des tickets. Le décrément n'a lieu que si le compteur est positif et le
lien actif ; le ticket n'est créé que si ce visiteur n'en a pas.

\textbf{Ce qui rend la faute possible, et comment elle est évitée.} Quatre
pièges ont été identifiés :

\begin{enumerate}[leftmargin=1.4em,itemsep=2pt]
  \item chercher le ticket existant \emph{après} le court-circuit « épuisé » :
        un visiteur déjà servi recevrait alors un refus. L'ordre est inversé
        exprès ;
  \item compter l'exemplaire au téléchargement du fichier et non à la
        réservation : un visiteur qui recharge la page consommerait deux fois.
        L'exemplaire est consommé au claim ;
  \item poser la transition « épuisé » dans la même écriture que le décrément :
        cinquante claims simultanés tenteraient alors cinquante fois la même
        transition. Elle est séparée et conditionnelle ;
  \item croire que la garde en mémoire décide du quota. Elle ne fait
        qu'\emph{alléger} la base sous le pic ; un visiteur qu'elle écarte reçoit
        « patienter », jamais « épuisé ».
\end{enumerate}

\textbf{Preuve.} Onze tests unitaires, dont un test de concurrence : cinquante
claims simultanés pour vingt exemplaires donnent exactement vingt réservations.
En charge, 15\,000 visiteurs sur un lien de 500 exemplaires ont produit
\emph{exactement} 500 réservations, sans une seule erreur serveur.

\subsubsection{C2 — aucune modification acceptée perdue}

\textbf{Où.} Dans l'écrivain unique par document et dans l'ordre
acquittement-après-écriture. Un lot est écrit en une transaction atomique, puis
diffusé en un seul message. Un acquittement ne part qu'après validation de la
transaction.

\textbf{Ce qui rend la faute possible, et comment elle est évitée.} Trois
risques :

\begin{enumerate}[leftmargin=1.4em,itemsep=2pt]
  \item acquitter avant d'écrire : une panne entre les deux perd la modification
        sans que personne ne le sache. L'acquittement est émis après ;
  \item renvoyer une modification sous un nouvel identifiant après une coupure :
        elle serait écrite deux fois. Le client renvoie le \emph{même}
        identifiant, et le serveur déduplique aussi par empreinte du contenu ;
  \item mesurer la convergence avec un outil qui ne peut pas la mesurer. Un
        générateur de charge compte des acquittements, il ne fait pas tourner un
        client CRDT.
\end{enumerate}

\textbf{Preuve.} Deux preuves distinctes, et c'est le point important. Le
générateur de charge vérifie que chaque modification émise reçoit son
acquittement. Un second outil, écrit en Node avec un vrai client CRDT,
reconstruit le document depuis le journal, vérifie que les numéros de séquence
sont contigus, et compare deux reconstructions indépendantes par empreinte. Ce
second outil a été validé dans les deux sens : il conclut à la convergence sur un
journal sain, et à l'échec dès qu'un numéro manque. Un détecteur qu'on n'a jamais
vu détecter ne prouve rien.

\subsubsection{C3 — document tel qu'il est à $T_1$}

\textbf{Où.} Au moment du claim : le dernier numéro de séquence réellement écrit
est figé dans le ticket. Le PDF est ensuite construit depuis le journal jusqu'à
ce numéro.

\textbf{Ce qui rend la faute possible.} Le PDF pourrait être généré à la demande,
au moment du téléchargement : deux visiteurs d'une même rafale recevraient alors
des documents différents, ce qui viole l'exigence de complétude et de cohérence.
Il pourrait aussi être régénéré à chaque visite, ce qui multiplierait le travail
par le nombre de visiteurs au lieu du nombre de versions.

\textbf{Preuve.} Le contenu du PDF est déterminé par le numéro figé, et la clé de
cache inclut ce numéro : une version, un fichier. Un verrou avec bail empêche
deux générations simultanées pour la même version.

\subsubsection{C4 — jamais plus que son crédit}

\textbf{Où.} Dans un compteur par adresse et par jour, débité dans la même
transaction que l'écriture au journal. Le comptage est effectué sur le binaire
Yjs, côté serveur : le client ne déclare pas ce qu'il a écrit.

\textbf{Ce qui rend la faute possible, et comment elle est évitée.}

\begin{enumerate}[leftmargin=1.4em,itemsep=2pt]
  \item compter côté client : un client modifié déclarerait zéro. Le comptage
        est serveur, en décodant la mise à jour ;
  \item ne pas facturer les suppressions : écrire puis effacer serait gratuit, et
        le plafond contournable en trois gestes. Les suppressions sont facturées ;
  \item compter par connexion plutôt que par adresse : cinq onglets donneraient
        cinq quotas. La clé est l'adresse ;
  \item inscrire le plafond comme condition de la transaction : DynamoDB refuse
        l'arithmétique entre valeurs. C'est une lecture préalable.
\end{enumerate}

\textbf{Preuve.} Seize tests unitaires, dont l'acceptation exacte jusqu'au
plafond et pas un caractère de plus, le partage du compteur entre plusieurs
connexions, la déduplication par identifiant et par empreinte, et le blocage au
quatrième refus.

\subsubsection{C5 — 15\,000 clics et 300 rédacteurs}

\textbf{Où.} À quatre endroits complémentaires : le micro-cache du serveur
frontal, la garde en mémoire sur les réservations, la diffusion par lots, et le
dimensionnement de l'instance.

\textbf{Ce qui rend la faute possible.} Chaque mécanisme a un mode de défaillance
propre, et c'est leur combinaison qui tient :

\begin{itemize}[leftmargin=1.4em,itemsep=2pt]
  \item sans micro-cache, 15\,000 lectures d'état frappent la base et la font
        échouer en limitation de débit ;
  \item sans garde en mémoire, autant d'écritures conditionnelles simultanées
        produisent le même effet ;
  \item sans diffusion par lots, un message par frappe multiplie la charge
        réseau par le nombre de rédacteurs ;
  \item sans dimensionnement mesuré, une instance trop petite s'effondre
        précisément pendant le pic qu'on mesure.
\end{itemize}

\textbf{Preuve.} Le scénario mixte fait tourner visiteurs, rédacteurs et sondes
d'état en même temps. Les mesures relevées sont la latence de réservation, la
latence d'état, la latence d'acquittement, et le nombre d'erreurs serveur —
celui-ci doit rester à zéro, y compris pour le parcours visiteur.

Un piège de mesure mérite d'être signalé : une réponse « épuisé » est un code 410,
que les outils de charge comptent comme une requête échouée. Un seuil sur ce taux
d'échec est donc inatteignable par construction — 84\,\% mesurés sans la moindre
panne. Le critère utile est un compteur dédié aux erreurs serveur.

\subsubsection{C6 — dernier exemplaire : prévenir vite, arrêter le travail}

\textbf{Où.} Dans la monotonie de l'état épuisé et dans son cache. Une fois le
lien épuisé, la réponse est immédiate, ne touche pas la base, et le serveur
frontal la conserve longuement. Le créateur du lien reçoit un courriel unique.

\textbf{Ce qui rend la faute possible.} Deux choses : recalculer l'état à chaque
visite, ce qui revient à travailler pour ne rien dire ; et envoyer le courriel
plusieurs fois, une par visiteur qui constate l'épuisement. Le courriel est donc
protégé par un verrou conditionnel : seul le claim qui pose la transition
déclenche l'envoi.

\textbf{Preuve.} L'état épuisé est posé une seule fois, et la notification est
idempotente : neuf tests unitaires couvrent l'envoi au seul dernier exemplaire,
l'absence de second envoi, la panne du service de courriel sans impact pour le
visiteur, et le relâchement du verrou en cas d'échec — mieux vaut un doublon rare
qu'un créateur jamais prévenu.

\begin{table}[H]
\centering
\small
\begin{tabular}{@{}p{0.7cm}p{5.6cm}p{8.7cm}@{}}
\toprule
\textbf{\#} & \textbf{Mécanisme} & \textbf{Preuve exécutable} \\
\midrule
C1 & transaction conditionnelle, ordre des vérifications, transition séparée
   & 11 tests dont concurrence 50/20 ; en charge, 500 réservations exactes \\
\addlinespace
C2 & écrivain unique, acquittement après écriture, déduplication par empreinte
   & acquittements sous charge ; contiguïté du journal et double reconstruction \\
\addlinespace
C3 & numéro figé au claim, une génération par version, verrou avec bail
   & clé de cache par version ; un seul travail par version \\
\addlinespace
C4 & comptage serveur sur le binaire, suppressions facturées, clé par adresse
   & 16 tests : plafond exact, multi-connexions, blocage au 4\up{e} refus \\
\addlinespace
C5 & micro-cache, garde mémoire, diffusion par lots, dimensionnement mesuré
   & scénario mixte : visiteurs, rédacteurs et sondes simultanés \\
\addlinespace
C6 & état épuisé monotone et mis en cache, notification idempotente
   & 9 tests : envoi unique, panne sans impact, verrou relâché \\
\bottomrule
\end{tabular}
\caption{Synthèse : une contrainte, un mécanisme, une preuve.}
\label{tab:synthese-contraintes}
\end{table}

% =============================================================================
% Partie C — Conclusion
% =============================================================================
\section{Conclusion}

\subsection{Ce qui est tenu, et comment on le sait}

Deux contraintes étaient éliminatoires. Elles sont traitées par des mécanismes
différents, et vérifiées par des preuves différentes — ce qui est précisément
pourquoi elles méritent d'être distinguées.

\textbf{C1 est prouvée en charge, en production.} Un lien de 500 exemplaires a
reçu 16\,753 requêtes dont 500 réservations exactement, sans une seule erreur
serveur, avec un pic à 200 requêtes par seconde. Le compteur final est à zéro et
l'état du lien est « épuisé ». Ce résultat n'est pas une moyenne : la contrainte
est « exactement $N$ », et l'écart mesuré est nul.

\textbf{C2 est prouvée en deux temps.} Sous charge, chaque modification émise a
reçu son acquittement : aucun message manquant. Mais un acquittement ne prouve
pas la convergence du contenu, et un outil de charge ne peut pas le faire. Un
second outil, avec un vrai client CRDT, reconstruit le document depuis le journal,
vérifie que les numéros de séquence sont contigus, et compare deux
reconstructions indépendantes. Il conclut à la convergence, et il conclut à
l'échec dès qu'un numéro manque.

\textbf{Les quatre autres contraintes sont couvertes par des tests unitaires.}
Le quota de crédit, la version figée du PDF, le comportement sous pic et la
notification du dernier exemplaire ont chacun leur jeu de tests, sans aucun appel
à un service distant.

\subsection{Ce qui reste à mesurer}

Le dossier doit être honnête sur ce point.

Le nombre de rédacteurs simultanés effectivement tenu reste à établir par paliers
successifs, et le résultat doit être rapporté tel quel, même inférieur à la cible
de 300. Une cible non atteinte et annoncée vaut mieux qu'un chiffre flatteur et
invérifiable.

Le comportement après un arrêt brutal du processus en pleine écriture dispose
d'un scénario dédié, qui tue le processus et vérifie que le journal a progressé
sans trou après redémarrage. Ce scénario est écrit, son exécution reste à
consigner.

Les métriques système — occupation processeur, mémoire de pointe, solde de
crédits de processeur — ne sont pas observables depuis l'outil de charge. Elles
se relèvent en parallèle sur la supervision du fournisseur, et doivent
accompagner les résultats de charge pour que le dimensionnement soit justifié
plutôt qu'affirmé.

\subsection{Ce que ce projet a réellement appris}

Trois leçons, dans l'ordre de ce qu'elles ont coûté.

\textbf{Première leçon : un faux service doit imiter les refus, pas seulement les
succès.} Un double en mémoire de la base évaluait lui-même les expressions de
condition. Il était donc plus permissif que le vrai service. Résultat :
quarante-deux tests verts, et aucune écriture possible en production, parce que le
service réel refuse l'arithmétique entre deux valeurs dans une condition. Le même
défaut s'est reproduit avec les mots réservés. La correction n'a pas été de
réparer le code, mais de rendre le double \emph{plus strict} : il rejette
désormais ce que le service rejette.

\textbf{Deuxième leçon : une compilation réussie ne prouve rien.} Les douze
diagrammes de conception compilaient sans erreur, et plusieurs étaient faux : un
code 200 là où le service répond 202, une redirection qui n'existe pas, un
numéro de version compté deux fois. L'outil de diagrammes écrit ses erreurs dans
l'image et renvoie un succès. Il a fallu relire chaque source contre le code et
regarder chaque rendu.

\textbf{Troisième leçon : un test qui ne peut pas échouer ne protège rien.} Un
seuil de taux d'échec posé sur un parcours qui renvoie légitimement des refus est
inatteignable par construction. Un outil de vérification qui interroge une route
dont la réponse n'a pas la forme attendue ne vérifie rien. À chaque fois, la
question utile n'est pas « le test passe-t-il ? » mais « ai-je déjà vu ce test
échouer ? ». C'est pour cette raison que les trois détecteurs de ce projet ont
été éprouvés en provoquant la panne qu'ils prétendent détecter.

\subsection{Perspectives}

Trois évolutions seraient utiles sans être nécessaires.

Le chiffrement du trafic suppose un nom de domaine : c'est la limite la plus
visible du déploiement actuel, et la plus simple à lever.

La haute disponibilité suppose de traiter l'écrivain unique autrement :
élection d'un responsable par document, ou externalisation de l'état. C'est un
chantier d'architecture à part entière, pas un réglage.

Le palier de rédacteurs, enfin, se mesure avant de s'optimiser : le résultat
réel doit guider la suite, plutôt qu'une optimisation supposée.

\subsection{Mot de la fin}

Le sujet prévenait : une application simple dont chaque choix est compris vaut
mieux qu'une application complexe qu'on ne sait pas expliquer. Ce dossier a
suivi cette consigne à la lettre, y compris dans ce qu'elle a de désagréable :
les limites sont énoncées, les hypothèses sont marquées comme telles, et les
écarts entre le plan et l'implémentation sont rapportés plutôt que lissés. C'est
la partie du travail qui a le plus servi, parce que c'est celle qui a révélé les
défauts que les tests ne voyaient pas.


\end{document}
